%% ============================================================
%%  fisher.tex  (revised)
%%  Fisher Information: From Classical Statistics to
%%  Quantum Neural Networks
%%
%%  Revised to align with week15_merged.tex:
%%    Slide 5  -- QFI Matrix, Fubini-Study, covariance form,
%%               circuit formula, parameter-shift connection
%%    Slide 13 -- Natural Gradient and TDVP
%%    Slide 14 -- Effective Dimension and Expressivity
%%    Slide 15 -- Barren Plateaus: Definition and Origin
%%
%%  Notation and macros follow week15_merged.tex exactly.
%% ============================================================

\documentclass[aspectratio=169,10pt]{beamer}

%% ── Theme & colours ──────────────────────────────────────────
\usetheme{Madrid}
\usecolortheme{default}
\usefonttheme[onlymath]{serif}

\setbeamercolor{frametitle}{bg=blue!12!white,fg=blue!60!black}
\setbeamercolor{title}{fg=white}
\setbeamercolor{block title}{bg=blue!18!white,fg=blue!60!black}
\setbeamercolor{block body}{bg=blue!4!white}
\setbeamercolor{alerted text}{fg=red!65!black}
\setbeamertemplate{navigation symbols}{}
\setbeamertemplate{footline}{%
  \leavevmode\hbox{%
    \begin{beamercolorbox}[wd=.25\paperwidth,ht=2.25ex,dp=1ex,center]
        {author in head/foot}\usebeamerfont{author in head/foot}
        \insertshortauthor
    \end{beamercolorbox}%
    \begin{beamercolorbox}[wd=.50\paperwidth,ht=2.25ex,dp=1ex,center]
        {title in head/foot}\usebeamerfont{title in head/foot}
        \insertshorttitle
    \end{beamercolorbox}%
    \begin{beamercolorbox}[wd=.25\paperwidth,ht=2.25ex,dp=1ex,right]
        {date in head/foot}\usebeamerfont{date in head/foot}
        \insertframenumber{}/\inserttotalframenumber\hspace*{1ex}
    \end{beamercolorbox}}%
  \vskip0pt}

%% ── Packages ─────────────────────────────────────────────────
\usepackage[utf8]{inputenc}
\usepackage[T1]{fontenc}
\usepackage{amsmath,amssymb,amsfonts,mathtools,bm}
\usepackage{booktabs}
\usepackage{xcolor}
\usepackage{hyperref}

%% ── Macros (identical to week15_merged.tex) ──────────────────
\newcommand{\btheta}{\boldsymbol{\theta}}
\newcommand{\bTheta}{\boldsymbol{\Theta}}
\newcommand{\QFI}{\mathbf{F}}
\newcommand{\E}{\mathbb{E}}
\newcommand{\Variance}{\mathrm{Var}}
\newcommand{\manifold}{\mathcal{M}}
\newcommand{\id}{\mathbb{1}}
\newcommand{\Tr}{\operatorname{Tr}}
\newcommand{\half}{\tfrac{1}{2}}
\newcommand{\R}{\mathbb{R}}
\newcommand{\C}{\mathbb{C}}
\newcommand{\hilb}{\mathcal{H}}
\newcommand{\FI}{\mathcal{I}}    % classical FI (scalar/matrix)

%% ── Metadata ──────────────────────────────────────────────────
\title[Fisher Information and QNNs]{%
  Fisher Information: From Classical Statistics\\
  to Quantum Neural Networks}
\subtitle{Mathematical Foundations, the QFI Matrix,
  and Optimisation Geometry}
\author[Graduate Lecture Notes]{Morten Hjorth-Jensen}
\institute{Department of Physics, University of Oslo, Norway}
\date{May 6, 2026}

%% ============================================================
\begin{document}
%% ============================================================

\begin{frame}[plain]\titlepage\end{frame}

\begin{frame}[allowframebreaks]{Outline}
  \scriptsize\tableofcontents[hidesubsections]
\end{frame}

%% ============================================================
\section{Classical Fisher Information}
%% ============================================================

\begin{frame}{Statistical Model and the Score Function}
  Parametric probability distribution $p(x|\theta)$,
  $\theta\in\R^d$.

  \medskip
  \textbf{Score function:}
  $s(x;\theta) = \partial_\theta \log p(x|\theta).$

  Key identity from $\int p\,dx=1$:
  $\E_{x\sim p}[s(x;\theta)] = 0.$
  The score has zero mean; its variance is the Fisher information.

  \medskip
  \begin{columns}[T]
  \column{0.52\textwidth}
    \textbf{Goals}
    \begin{itemize}\itemsep2pt
      \item Estimate $\theta$ from data
      \item Quantify sensitivity of $p$ to $\theta$
      \item Define a geometry on parameter space
    \end{itemize}
  \column{0.45\textwidth}
    \begin{block}{Intuition}
      \footnotesize
      High sensitivity $\Rightarrow$ nearby distributions
      are distinguishable $\Rightarrow$ more information
      about $\theta$ in a single observation.
    \end{block}
  \end{columns}
\end{frame}

%% ─────────────────────────────────────────────────────────────
\begin{frame}{Definition and Equivalent Forms}
  \textbf{Fisher Information} (scalar $\theta$):
  \[
    \FI(\theta)
    = \E_{x\sim p}\!\left[
        \bigl(\partial_\theta \log p(x|\theta)\bigr)^2
      \right]
    = \Variance[s(x;\theta)].
  \]

  \textbf{Equivalent form:}
  \[
    \FI(\theta) = -\,\E\!\left[\partial_\theta^2 \log p(x|\theta)\right].
  \]

  \textbf{Matrix version} for $\theta\in\R^d$:
  \[
    \FI_{jk}(\theta)
    = \E\!\left[
        \frac{\partial\log p}{\partial\theta_j}
        \frac{\partial\log p}{\partial\theta_k}
      \right]
    = -\,\E\!\left[
        \frac{\partial^2\log p}{\partial\theta_j\partial\theta_k}
      \right].
  \]

  \begin{block}{Derivation sketch}
    Differentiate $\E[\partial_\theta\log p]=0$ once more:
    $\E[\partial_\theta^2\log p]+\E[(\partial_\theta\log p)^2]=0$,
    giving the equivalence immediately.
  \end{block}
\end{frame}

%% ─────────────────────────────────────────────────────────────
\begin{frame}{Cramér--Rao Bound and Natural Gradient}
  \begin{columns}[T]
  \column{0.52\textwidth}
    \textbf{Cramér--Rao bound:}
    for any unbiased $\hat\theta$,
    \[
      \Variance(\hat\theta) \geq \frac{1}{\FI(\theta)}.
    \]
    Matrix version: $\mathrm{Cov}(\hat\theta)\geq\FI(\theta)^{-1}$
    (Löwner order).
    \begin{itemize}\itemsep2pt
      \item $\FI$ sets a fundamental precision limit
      \item High $\FI$ $\Rightarrow$ sharper inference
      \item MLEs saturate the bound asymptotically
    \end{itemize}
  \column{0.45\textwidth}
    \textbf{Natural gradient} (Amari 1998):
    steepest descent on the manifold $(\R^d,\FI)$:
    \[
      \theta_{t+1}
      = \theta_t - \eta\,\FI(\theta)^{-1}\,\nabla_\theta L.
    \]
    \begin{block}{Fisher--Rao metric}
      \[
        ds^2 = \sum_{jk}\FI_{jk}\,d\theta_j\,d\theta_k.
      \]
      Reparameterisation-invariant;
      steepest descent on information-geometric manifold.
    \end{block}
  \end{columns}
\end{frame}

%% ============================================================
\section{Quantum Fisher Information}
%% ============================================================

\begin{frame}{From Classical to Quantum}
  \begin{columns}[T]
  \column{0.52\textwidth}
    Classical: $p(x|\theta) \to$ density matrix
    $\rho(\btheta)$.

    \medskip
    \textbf{Pure variational state:}
    \[
      |\psi(\btheta)\rangle = U(\btheta)|0\rangle^{\otimes n},
      \quad
      \rho = |\psi\rangle\langle\psi|.
    \]
    \textbf{Observable expectation:}
    \[
      f(\btheta)
      = \langle\psi(\btheta)|\,\hat{O}\,|\psi(\btheta)\rangle.
    \]
    \textbf{SLD} (Symmetric Logarithmic Derivative):
    \[
      \partial_\theta\rho
      = \frac{1}{2}(L_\theta\rho + \rho L_\theta).
    \]
    \textbf{QFI (mixed state):}
    $F_Q = \Tr[\rho\,L_\theta^2].$
  \column{0.45\textwidth}
    \begin{block}{Quantum Cramér--Rao bound}
      For any measurement strategy:
      \[
        \Variance(\hat\theta)\geq \frac{1}{F_Q(\theta)}.
      \]
      $F_Q\geq\FI_{\rm classical}$ for every measurement.
      The QFI maximises over all quantum measurements.
    \end{block}
    \begin{alertblock}{Key difference}
      The classical FI depends on a measurement choice.
      The QFI is the \emph{maximum} over all possible
      quantum measurements---the ultimate precision limit.
    \end{alertblock}
  \end{columns}
\end{frame}

%% ─────────────────────────────────────────────────────────────
\begin{frame}{QFI for Pure States and the Generator Formula}
  For $|\psi(\theta)\rangle$ (single parameter):
  \[
    F_Q(\theta)
    = 4\Bigl(
        \langle\partial_\theta\psi|\partial_\theta\psi\rangle
       -\bigl|\langle\psi|\partial_\theta\psi\rangle\bigr|^2
      \Bigr).
  \]

  \textbf{Generator form.}
  If $|\psi(\theta)\rangle = e^{-i\theta G}|\psi_0\rangle$
  with Hermitian generator $G$ ($G^\dagger=G$), then
  $\partial_\theta|\psi\rangle = -iG|\psi\rangle$ and:
  \[
    F_Q = 4\,\Variance_{|\psi\rangle}(G)
        = 4\bigl(\langle G^2\rangle - \langle G\rangle^2\bigr).
  \]

  \begin{columns}[T]
  \column{0.52\textwidth}
    \textbf{Example: single qubit.}
    $|\psi(\theta)\rangle = e^{-i\theta\sigma_y/2}|0\rangle$,
    $G=\sigma_y/2$, eigenvalues $\pm\tfrac12$:
    \[
      F_Q = 4\,\Variance(\sigma_y/2)
          = 4\cdot\tfrac14 = 1.
    \]
    Uniform sensitivity; no barren plateau.

    \textbf{Metrological scaling:}
    for $N$ separable qubits $F_Q\leq N$ (SQL);
    with entanglement $F_Q\sim N^2$ (Heisenberg limit).
  \column{0.45\textwidth}
    \begin{alertblock}{Interpretation}
      $F_Q = 4\,\Variance(G)$: large generator fluctuations
      mean the state is highly sensitive to $\theta$;
      information about $\theta$ is efficiently encoded.
      Maximising $\Variance(G)$ is the task of
      quantum metrology.
    \end{alertblock}
  \end{columns}
\end{frame}

%% ============================================================
\section{The QFI Matrix in QNNs}
%% ============================================================

%% ─── Aligns with week15.tex Slide 5 ─────────────────────────
\begin{frame}{The QFI Matrix: Definition and Geometric Meaning}
\small
  For a variational state $|\psi(\btheta)\rangle$,
  $\btheta\in\R^p$, the QFI is a
  \emph{real symmetric $p\times p$ matrix}:
  \[
    \boxed{
      F_{jk}(\btheta)
      = 4\,\mathrm{Re}\!\Bigl(
          \langle\partial_j\psi|\partial_k\psi\rangle
         -\langle\partial_j\psi|\psi\rangle
          \langle\psi|\partial_k\psi\rangle
        \Bigr)
    }
  \]
  $|\partial_j\psi\rangle
    \equiv \partial|\psi(\btheta)\rangle/\partial\theta_j$.

  \begin{columns}[T]
  \column{0.52\textwidth}
    \textbf{Geometric meaning.}
    $\QFI$ is the \textbf{Fubini--Study metric}
    on the variational manifold $\manifold$:
    \[
      ds^2 = \sum_{jk} F_{jk}\,d\theta_j\,d\theta_k.
    \]
    Measures information-geometric distance between
    neighbouring quantum states in $\hilb$.

    \medskip
    \textbf{Covariance form}
    (product-of-rotations circuit
    $U(\btheta)=\prod_j e^{-i\theta_j P_j/2}$):
    \[
      F_{jk} = 4\,\mathrm{Cov}_{|\psi\rangle}(G_j, G_k),
      \quad G_j = P_j/2.
    \]
  \column{0.45\textwidth}
    \begin{block}{Circuit formula (two-fold shifts)}
      State-overlap function:
      $\mathcal{O}(\btheta)
        = |\langle\psi_{\rm ref}|\psi(\btheta)\rangle|^2$.
      \begin{align*}
        F_{jk} &= \tfrac12\bigl[
          \mathcal{O}_{j+,k+}
         -\mathcal{O}_{j+,k-}\\
         &\qquad -\mathcal{O}_{j-,k+}
         +\mathcal{O}_{j-,k-}\bigr],
      \end{align*}
      $\mathcal{O}_{j\pm,k\pm}$: shift
      $\theta_j\to\theta_j\pm\tfrac\pi2$,
      $\theta_k\to\theta_k\pm\tfrac\pi2$.
      Cost: $O(p^2)$ circuit evaluations.
    \end{block}
  \end{columns}
\end{frame}

%% ─────────────────────────────────────────────────────────────
\begin{frame}{Derivation of the Covariance Form}
  For $U(\btheta)=\prod_j e^{-i\theta_j P_j/2}$
  in the single-layer case:
  $\partial_j|\psi\rangle = -\tfrac{i}{2}P_j|\psi\rangle$.

  \medskip
  Compute each term in the QFI formula:
  \begin{align*}
    \langle\partial_j\psi|\partial_k\psi\rangle
    &= \tfrac{1}{4}\langle\psi|P_j P_k|\psi\rangle, \\
    \langle\partial_j\psi|\psi\rangle
    &= \tfrac{i}{2}\langle P_j\rangle, \quad
    \langle\psi|\partial_k\psi\rangle
    = -\tfrac{i}{2}\langle P_k\rangle.
  \end{align*}

  Insert and take the real part:
  \[
    F_{jk}
    = 4\,\mathrm{Re}\!\left(
        \tfrac{1}{4}\langle P_j P_k\rangle
       -\tfrac{i}{2}\langle P_j\rangle\cdot
        \tfrac{i}{2}\langle P_k\rangle
      \right)
    = \langle P_j P_k\rangle
      - \langle P_j\rangle\langle P_k\rangle
    = 4\,\mathrm{Cov}(G_j,G_k). \quad\square
  \]

  \begin{alertblock}{Multi-layer extension}
    For $L$-layer circuits, commutators between generators
    of different layers contribute.
    The two-fold shift circuit formula handles all layers
    exactly at cost $O(p^2)$.
  \end{alertblock}
\end{frame}

%% ─────────────────────────────────────────────────────────────
\begin{frame}{Parameter-Shift Rule and the QFI Connection}
  \textbf{Parameter-shift rule} for
  $f(\btheta)=\langle\psi(\btheta)|\hat{O}|\psi(\btheta)\rangle$
  with generator eigenvalues $\pm\tfrac12$:
  \[
    \frac{\partial f}{\partial\theta_j}
    = \frac{1}{2}\!\left[
        f\!\left(\theta_j+\tfrac\pi2\right)
       -f\!\left(\theta_j-\tfrac\pi2\right)
      \right].
  \]
  Exact, no approximation; 2 circuit evaluations per parameter.

  \medskip
  \textbf{Commutator form:}
  $\partial_{\theta_j}f
   = \langle\psi|i[G_j,\hat{O}]|\psi\rangle$.

  \medskip
  \textbf{Diagonal QFI entry:}
  $F_{jj} = 4\,\Variance_{|\psi\rangle}(G_j)$.

  \medskip
  \begin{block}{Cauchy--Schwarz bound (gradient--QFI)}
    \[
      \left|\frac{\partial f}{\partial\theta_j}\right|^2
      \leq F_{jj}\cdot\Variance_{|\psi\rangle}(\hat{O}).
    \]
    Small generator variance $\Rightarrow$ small $F_{jj}$
    $\Rightarrow$ all gradients of all observables are small
    $\Rightarrow$ barren plateau for parameter $\theta_j$.
  \end{block}
\end{frame}

%% ─────────────────────────────────────────────────────────────
\begin{frame}{Quantum Geometric Tensor}
  The \textbf{Quantum Geometric Tensor (QGT)}:
  \[
    \chi_{jk}
    = \langle\partial_j\psi|\partial_k\psi\rangle
     -\langle\partial_j\psi|\psi\rangle
      \langle\psi|\partial_k\psi\rangle.
  \]

  Decomposition: $\chi_{jk} = \frac{1}{4}F_{jk} + i\,\Omega_{jk}$,
  \begin{itemize}\itemsep2pt
    \item $F_{jk} = 4\,\mathrm{Re}(\chi_{jk})$:
          \textbf{QFI matrix} (Fubini--Study metric)
    \item $\Omega_{jk} = \mathrm{Im}(\chi_{jk})$:
          \textbf{Berry curvature} (geometric phase)
  \end{itemize}

  \begin{columns}[T]
  \column{0.52\textwidth}
    \textbf{Geometric significance.}
    $\chi$ is the pullback of the Fubini--Study form on
    $\mathbb{P}(\hilb)$ via $\btheta\mapsto|\psi(\btheta)\rangle$.
    It unifies information geometry (real part) and
    symplectic geometry (imaginary part).
  \column{0.45\textwidth}
    \begin{alertblock}{Connection to Berry phase}
      The Berry phase around a loop $\gamma$
      in parameter space is:
      \[
        \varphi_B = \oint_\gamma
          \mathrm{Im}\langle\psi|d|\psi\rangle
        = \iint_\Sigma \Omega_{jk}\,d\theta_j\wedge d\theta_k.
      \]
    \end{alertblock}
  \end{columns}
\end{frame}

%% ============================================================
\section{Natural Gradient and TDVP}
%% ============================================================

%% ─── Aligns with week15.tex Slide 13 ────────────────────────
\begin{frame}{Natural Gradient in Quantum Circuits}
  \begin{columns}[T]
  \column{0.52\textwidth}
    \textbf{Natural gradient} (Amari 1998; Stokes et al.\ 2020):
    \[
      \btheta \;\leftarrow\; \btheta
        - \eta\,\QFI(\btheta)^{+}\,\nabla_\theta C.
    \]
    $\QFI^+$: Moore--Penrose pseudoinverse
    (needed when $\QFI$ is rank-deficient).

    \medskip
    \textbf{Derivation.}
    Minimise $C(\btheta+\delta\btheta)$ subject to
    fixed Bures distance:
    \[
      \delta\btheta^\top\QFI\,\delta\btheta = \varepsilon^2.
    \]
    Lagrange multipliers give
    $\delta\btheta\propto-\QFI^{-1}\nabla C$.

    \medskip
    The update is \emph{reparameterisation-invariant}:
    it always descends the steepest direction on $\manifold$.
  \column{0.44\textwidth}
    \begin{block}{TDVP equations of motion}
      Dirac--Frenkel variational principle:
      project $(i\partial_t - H)|\psi\rangle = 0$
      onto $\mathcal{T}_{\btheta}\manifold
        = \mathrm{span}\{|\partial_j\psi\rangle\}$:
      \[
        \sum_k F_{jk}\,\dot{\theta}_k
        = 2\,\mathrm{Im}\langle\partial_j\psi|H|\psi\rangle.
      \]
      This is the Time-Dependent Variational Principle
      for VQAs.
    \end{block}
    \begin{alertblock}{Imaginary-time TDVP $=$ natural gradient VQE}
      Substituting $t\to-i\tau$:
      $\dot{\btheta} = -\QFI^{-1}\nabla_\theta E(\btheta)$.
      Equivalent to \textbf{Stochastic Reconfiguration} in VMC.
    \end{alertblock}
  \end{columns}
\end{frame}

%% ─────────────────────────────────────────────────────────────
\begin{frame}{Computing the Natural Gradient and SR Connection}
  \begin{columns}[T]
  \column{0.52\textwidth}
    \textbf{Exact cost}: $O(p^2)$ circuit evaluations plus
    $O(p^3)$ matrix inversion. Expensive for $p\sim10^3$.

    \medskip
    \textbf{Practical approximations:}
    \begin{enumerate}\itemsep2pt
      \item \textbf{Diagonal}: use only $F_{jj}$;
            $O(p)$ evaluations.
      \item \textbf{Block-diagonal}: group by layer;
            $O(p^2/L)$ cost.
      \item \textbf{Regularised}: replace $\QFI$ by
            $\QFI+\lambda I$.
    \end{enumerate}

    \medskip
    \textbf{SR connection (Sorella 1998):}
    \[
      \theta_j \leftarrow \theta_j
        - \eta\sum_k S_{jk}^{-1} g_k,
    \]
    with $S_{jk} = \langle\partial_j\psi|\partial_k\psi\rangle
                 - \langle\partial_j\psi|\psi\rangle
                   \langle\psi|\partial_k\psi\rangle
                 = F_{jk}/4$.
    SR is the natural gradient w.r.t.\ the Fubini--Study metric.
  \column{0.45\textwidth}
    \begin{block}{Regularised update rule}
      \[
        \btheta \leftarrow \btheta
          - \eta\,(\QFI + \lambda I)^{-1}\,\nabla_\theta C.
      \]
      \begin{itemize}\itemsep2pt
        \item $\lambda\to 0$: pure natural gradient
        \item $\lambda\to\infty$: standard gradient
        \item $\lambda = \|\nabla C\|$: adaptive
      \end{itemize}
    \end{block}
    \begin{alertblock}{Quantum advantage}
      The QFI matrix elements $F_{jk}$ can be measured
      on a quantum device via the two-fold shift formula,
      without classical density-matrix simulation.
    \end{alertblock}
  \end{columns}
\end{frame}

%% ============================================================
\section{Effective Dimension and Expressivity}
%% ============================================================

%% ─── Aligns with week15.tex Slide 14 ────────────────────────
\begin{frame}{Effective Dimension via the QFI}
  \textbf{Abbas et al.\ (2021)}: define the
  \textbf{effective dimension} via the average log-volume
  of the QFI:
  \[
    d_{\rm eff}
    = \frac{2}{\log N}
      \log\!\left[\frac{1}{V}\int_{\btheta}
        \sqrt{\det\!\left(
          I + \frac{N}{2\pi\ln N}\,\QFI(\btheta)
        \right)}\,d\btheta\right],
  \]
  where $N$ is the number of training samples.

  \begin{columns}[T]
  \column{0.52\textwidth}
    \textbf{Properties}
    \begin{itemize}\itemsep1pt
      \item High $d_{\rm eff}$: large effective capacity.
      \item Some QNNs exceed $d_{\rm eff}$ of classical
            networks at equal parameter count.
      \item $\QFI\approx 0$: redundant parameters;
            flat loss directions.
      \item $d_{\rm eff}\to 0$ as $\QFI\to 0$: model
            collapses.
      \item Maximal expressibility (Haar) gives high
            $d_{\rm eff}$ \emph{but also barren plateaus}.
    \end{itemize}
  \column{0.44\textwidth}
    \begin{block}{QFI as trainability indicator}
      Cauchy--Schwarz on the Fubini--Study metric:
      \[
        \left|\frac{\partial C}{\partial\theta_j}\right|^2
        \leq
        \lambda_{\max}(\QFI)
        \cdot
        \Variance_{|\psi\rangle}(\hat{B}).
      \]
      Small $\lambda_{\max}(\QFI)$
      $\Leftrightarrow$ all gradients small
      $\Leftrightarrow$ barren plateau.
      Monitor $\QFI$ spectrum as an early-warning diagnostic.
    \end{block}
  \end{columns}
\end{frame}

%% ─────────────────────────────────────────────────────────────
\begin{frame}{QFI Rank, Redundant Parameters, and Expressibility}
  \begin{columns}[T]
  \column{0.52\textwidth}
    \textbf{Rank deficiency.}
    If $\mathrm{rank}(\QFI)<p$, there exist directions
    $\mathbf{v}$ with $\mathbf{v}^\top\QFI\,\mathbf{v}=0$:
    moving along $\mathbf{v}$ does not change the quantum state
    $\Rightarrow$ \emph{redundant parameters}.

    \medskip
    \textbf{Consequences:}
    \begin{itemize}\itemsep2pt
      \item Flat loss directions (pseudo-barren)
      \item Can aid optimisation (escape local minima)
      \item Inflate circuit depth needlessly
      \item Use $\mathrm{rank}(\QFI)$ to prune circuits
    \end{itemize}

    \medskip
    \textbf{Expressibility} (Meyer et al.\ 2021):
    measures how uniformly the circuit covers $U(2^n)$.
    High expressibility $\Leftrightarrow$ large
    $\det(\QFI)$ on average $\Leftrightarrow$ high $d_{\rm eff}$.
  \column{0.45\textwidth}
    \begin{alertblock}{The expressibility--trainability trade-off}
      \begin{center}
      \begin{tabular}{lcc}
        \toprule
        Circuit & $d_{\rm eff}$ & Trainable? \\
        \midrule
        Shallow, structured & Low & Yes \\
        Medium depth & Medium & Often \\
        Deep random & High & No (BP) \\
        \bottomrule
      \end{tabular}
      \end{center}
      Optimal VQAs balance expressibility with
      structured ans\"{a}tze that avoid the
      Haar-random concentration of measure.
    \end{alertblock}
  \end{columns}
\end{frame}

%% ============================================================
\section{Barren Plateaus and the QFI}
%% ============================================================

%% ─── Aligns with week15.tex Slide 15 ────────────────────────
\begin{frame}{Barren Plateaus: Definition and Mathematical Origin}
\small
  \begin{block}{Definition (McClean et al.\ 2018)}
    A circuit family $\{U(\btheta)\}$ has a
    \textbf{barren plateau} if:
    \[
      \E_{\btheta}\!\left[\partial_{\theta_j}C\right] = 0,
      \qquad
      \Variance_{\btheta}\!\left[\partial_{\theta_j}C\right]
      = O(b^{-n}),\quad b>1.
    \]
  \end{block}

  \begin{columns}[T]
  \column{0.52\textwidth}
    \textbf{Mathematical origin.}
    For Haar-random $U$ on $\mathrm{SU}(2^n)$:
    \[
      \E_U[\langle\hat{O}\rangle]
      = \tfrac{\Tr\hat{O}}{2^n},\;
      \Variance[\partial_{\theta_j}C]
      \sim \tfrac{\Tr\hat{O}^2}{4^n}.
    \]
    For $\hat{O}=Z_k$: $\Variance\sim 1/2^n$.
    Concentration of measure (Lévy's lemma):
    functions concentrate near their mean with prob.\
    $\geq 1-2e^{-\Omega(2^{2n}\varepsilon^2)}$.
    \begin{alertblock}{Consequence}
      Deep random circuits:
      $C(\btheta)\approx\Tr\hat{O}/2^n$ almost everywhere.
      Flat landscape. \emph{Not} a local-minima problem.
    \end{alertblock}
  \column{0.44\textwidth}
    \begin{block}{Direct QFI link}
      Cauchy--Schwarz:
      $\Variance[\partial_{\theta_j}C]
        \leq \tfrac{1}{4}F_{jj}\Variance(\hat{O})$.

      For Haar-random circuits:
      $F_{jj}\sim 1/2^n\to 0$.

      Therefore:
      $F_{jj}\to 0 \Longleftrightarrow \text{barren plateau.}$

      Monitoring $F_{jj}$ before and during training
      gives an \emph{early warning} of trainability collapse.
    \end{block}
  \end{columns}
\end{frame}

%% ─────────────────────────────────────────────────────────────
\begin{frame}{Noise-Induced Barren Plateaus and Mitigation}
  \begin{columns}[T]
  \column{0.52\textwidth}
    \textbf{Hardware noise} (Wang et al.\ 2021):
    depolarising noise rate $\varepsilon$ per gate:
    \[
      |\langle\hat{O}\rangle_{\rm noisy}|
      \leq (1-\varepsilon)^{pd}\,
      |\langle\hat{O}\rangle_{\rm ideal}|,
    \]
    \[
      \Variance[\partial_{\theta_j}C_{\rm noisy}]
      \sim (1-\varepsilon)^{2pd}
           \,\Variance[\partial_{\theta_j}C_{\rm ideal}].
    \]
    Affects even \emph{shallow} circuits;
    cannot be fixed by initialisation.
  \column{0.44\textwidth}
    \begin{block}{Mitigation strategies}
      \footnotesize
      \begin{tabular}{ll}
        \toprule
        \textbf{Strategy} & \textbf{Mechanism} \\
        \midrule
        Local cost function & poly$(n)$ gradient scaling \\
        Problem ansatz & exploit symmetries \\
        Layer-wise training & $O(1)$ depth locally \\
        Identity initialisation & $O(1)$ initial gradients \\
        Natural gradient & QFI-respecting steps \\
        Shallow circuits & avoid Haar regime \\
        \bottomrule
      \end{tabular}
    \end{block}
    \begin{alertblock}{Physics-inspired ansätze}
      Circuits mirroring Hamiltonian structure (QPINNs)
      naturally avoid the Haar regime and mitigate BPs
      without additional overhead.
    \end{alertblock}
  \end{columns}
\end{frame}

%% ─────────────────────────────────────────────────────────────
\begin{frame}{Entanglement, QNTK, and Training Regimes}
  \begin{columns}[T]
  \column{0.52\textwidth}
    \textbf{Entanglement entropy and trainability.}
    $S(\rho_A)=-\Tr(\rho_A\log\rho_A)$:

    \medskip
    \begin{center}
    \begin{tabular}{lll}
      \toprule
      Regime & $S(\rho_A)$ & Training \\
      \midrule
      Low & $\approx 0$ & Classically simulable \\
      Volume-law & $\Theta(n)$ & Barren plateaus \\
      Area-law & $O(|\partial A|)$ & Trainable \\
      \bottomrule
    \end{tabular}
    \end{center}
    Area-law (shallow circuits, MERA) is the sweet spot
    for VQA trainability.
  \column{0.44\textwidth}
    \begin{block}{Quantum Neural Tangent Kernel (QNTK)}
      \[
        \Theta(\mathbf{x},\mathbf{x}')
        = \sum_j
          \frac{\partial f(\mathbf{x};\btheta)}{\partial\theta_j}
          \frac{\partial f(\mathbf{x}';\btheta)}{\partial\theta_j}.
      \]
      \begin{itemize}\itemsep2pt
        \item Governs linearised training dynamics
        \item At a barren plateau:
              $\Theta\approx 0$ for all $\mathbf{x},\mathbf{x}'$
        \item Bound: $\Theta(\mathbf{x},\mathbf{x})
              \leq\lambda_{\max}(\QFI)$
        \item Connects QNN training to quantum kernel SVMs
      \end{itemize}
    \end{block}
  \end{columns}
\end{frame}

%% ============================================================
\section{Summary}
%% ============================================================

\begin{frame}{Key Equations at a Glance}
  \small
  \begin{center}
  \renewcommand{\arraystretch}{1.2}
  \begin{tabular}{p{3.8cm}p{7.2cm}}
    \toprule
    \textbf{Concept} & \textbf{Formula} \\
    \midrule
    Classical FI &
      $\FI(\theta)=\E[(\partial_\theta\log p)^2]
                 =-\E[\partial_\theta^2\log p]$ \\
    Cramér--Rao &
      $\Variance(\hat\theta)\geq 1/\FI(\theta)$ \\
    Classical nat.\ grad. &
      $\theta\leftarrow\theta-\eta\,\FI^{-1}\nabla L$ \\
    QFI matrix &
      $F_{jk}=4\,\mathrm{Re}(\langle\partial_j\psi|\partial_k\psi\rangle
             -\langle\partial_j\psi|\psi\rangle
              \langle\psi|\partial_k\psi\rangle)$ \\
    Covariance form &
      $F_{jk}=4\,\mathrm{Cov}_{|\psi\rangle}(G_j,G_k)$ \\
    Generator variance &
      $F_Q=4\,\Variance_{|\psi\rangle}(G)$ \\
    Fubini--Study metric &
      $ds^2=\sum_{jk}F_{jk}\,d\theta_j d\theta_k$ \\
    Quantum geometric tensor &
      $\chi_{jk}=\frac{1}{4}F_{jk}+i\Omega_{jk}$ \\
    Parameter-shift &
      $\partial_{\theta_j}f=\frac{1}{2}[f(\theta_j+\frac\pi2)
                                       -f(\theta_j-\frac\pi2)]$ \\
    Gradient--QFI bound &
      $|\partial_{\theta_j}C|^2\leq\lambda_{\max}(\QFI)\cdot\Variance(\hat{O})$ \\
    Natural grad.\ (quantum) &
      $\btheta\leftarrow\btheta-\eta\,\QFI^+\nabla C$ \\
    TDVP &
      $\sum_k F_{jk}\dot\theta_k=2\,\mathrm{Im}\langle\partial_j\psi|H|\psi\rangle$ \\
    Barren plateau &
      $\Variance[\partial_{\theta_j}C]\sim 1/2^n$,
      $F_{jj}\sim 1/2^n$ \\
    \bottomrule
  \end{tabular}
  \end{center}
\end{frame}

%% ─────────────────────────────────────────────────────────────
\begin{frame}{Summary}
  \begin{columns}[T]
  \column{0.52\textwidth}
    \textbf{Classical Fisher information}
    \begin{itemize}\itemsep2pt
      \item Variance of the score function
      \item Cramér--Rao precision limit
      \item Fisher--Rao Riemannian metric
      \item Natural gradient: reparameterisation-invariant
            descent
    \end{itemize}
    \textbf{Quantum generalisation}
    \begin{itemize}\itemsep2pt
      \item QFI $=$ maximum FI over all measurements
      \item Pure-state formula via state overlaps
      \item QFI matrix $=$ Fubini--Study metric on $\manifold$
      \item Covariance of generators; exact circuit formula
      \item QGT: QFI (metric) + Berry curvature (symplectic)
    \end{itemize}
  \column{0.45\textwidth}
    \textbf{QFI in QNN optimisation}
    \begin{itemize}\itemsep2pt
      \item Natural gradient $= \QFI^{-1}\nabla C$:
            geometry-aware, invariant
      \item TDVP: imaginary-time $=$ natural gradient VQE
            $=$ stochastic reconfiguration
      \item Effective dimension via $\det(\QFI)$:
            quantifies model capacity
      \item Barren plateaus: $F_{jj}\sim 1/2^n$;
            QFI spectrum is an early-warning diagnostic
    \end{itemize}
    \begin{alertblock}{Central message}
      The QFI is simultaneously the
      \emph{optimisation metric},
      the \emph{expressivity measure}, and the
      \emph{barren-plateau detector}
      for quantum neural networks.
    \end{alertblock}
  \end{columns}
\end{frame}

%% ─────────────────────────────────────────────────────────────
\begin{frame}{Exercises}
  \begin{enumerate}\itemsep3pt
    \item \textbf{Score identity.}
          Verify $\E[\partial_\theta\log p]=0$ for
          $p(x|\theta)=\mathcal{N}(x;\theta,1)$;
          compute $\FI(\theta)$.

    \item \textbf{QFI of a single qubit.}
          For $|\psi(\theta)\rangle=R_y(\theta)|0\rangle$,
          compute $F_{11}$ via the Fubini--Study formula
          and via $4\,\Variance(G)$.
          Verify $F_{11}=1$.

    \item \textbf{Two-parameter covariance form.}
          For $U(\theta_1,\theta_2)
            =e^{-i\theta_1 Z/2}e^{-i\theta_2 X/2}$,
          compute the full $2\times 2$ matrix
          $F_{jk}=4\,\mathrm{Cov}(G_j,G_k)$.

    \item \textbf{Circuit formula verification.}
          Verify the two-fold shift formula for $F_{12}$
          numerically for the above circuit on $|0\rangle$.

    \item \textbf{Barren plateau scaling.}
          Estimate
          $\Variance[\partial C/\partial\theta_j]$ for
          $n=2,\ldots,8$ qubit random circuits of depth $2n$.
          Compare with diagonal $F_{jj}$.

    \item \textbf{Natural gradient vs.\ Adam.}
          Implement both for a two-qubit VQE problem;
          compare convergence curves.
  \end{enumerate}
\end{frame}

%% ─────────────────────────────────────────────────────────────
\begin{frame}[allowframebreaks]{References}
  \begin{columns}[T]
  \column{0.50\textwidth}
    \textbf{Classical Fisher information}
    \begin{enumerate}\scriptsize
      \item S.\ Amari,
            \emph{Natural gradient works efficiently in learning},
            Neural Comput.\ \textbf{10}, 251 (1998).
      \item S.\ Amari \& H.\ Nagaoka,
            \emph{Methods of Information Geometry},
            AMS/Oxford (2000).
    \end{enumerate}
    \textbf{Quantum Fisher information}
    \begin{enumerate}\scriptsize\setcounter{enumi}{2}
      \item S.\ L.\ Braunstein \& C.\ M.\ Caves,
            \emph{Statistical distance and geometry of
            quantum states},
            PRL \textbf{72}, 3439 (1994).
      \item M.\ G.\ A.\ Paris,
            \emph{Quantum estimation for quantum technology},
            Int.\ J.\ Quantum Inf.\ \textbf{7}, 125 (2009).
    \end{enumerate}
    \textbf{QFI in variational algorithms}
    \begin{enumerate}\scriptsize\setcounter{enumi}{4}
      \item J.\ Stokes et al.,
            \emph{Quantum natural gradient},
            Quantum \textbf{4}, 269 (2020).
      \item A.\ Abbas et al.,
            \emph{The power of quantum neural networks},
            Nat.\ Comput.\ Sci.\ \textbf{1}, 403 (2021).
      \item S.\ Sorella,
            \emph{Stochastic reconfiguration},
            PRL \textbf{80}, 4558 (1998).
    \end{enumerate}
  \column{0.46\textwidth}
    \textbf{Barren plateaus}
    \begin{enumerate}\scriptsize\setcounter{enumi}{7}
      \item J.\ R.\ McClean et al.,
            \emph{Barren plateaus in QNN training landscapes},
            Nat.\ Commun.\ \textbf{9}, 4812 (2018).
      \item M.\ Cerezo et al.,
            \emph{Cost-function-dependent barren plateaus},
            Nat.\ Commun.\ \textbf{12}, 1791 (2021).
      \item S.\ Wang et al.,
            \emph{Noise-induced barren plateaus},
            Nat.\ Commun.\ \textbf{12}, 6961 (2021).
    \end{enumerate}
    \textbf{QNN expressivity and geometry}
    \begin{enumerate}\scriptsize\setcounter{enumi}{10}
      \item M.\ Cerezo et al.,
            \emph{Variational quantum algorithms},
            Nat.\ Rev.\ Phys.\ \textbf{3}, 625 (2021).
      \item A.\ P\'erez-Salinas et al.,
            \emph{Data re-uploading for a universal
            quantum classifier},
            Quantum \textbf{4}, 226 (2020).
      \item M.\ Schuld et al.,
            \emph{Effect of data encoding on expressivity},
            PRA \textbf{103}, 032430 (2021).
    \end{enumerate}
  \end{columns}
\end{frame}

\end{document}
